\documentclass[aspectratio=169]{beamer}
\usetheme{Madrid}
\usefonttheme{structurebold}
\usefonttheme[onlymath]{serif}

\definecolor{EricTemplate}{RGB}{255,103,0} 
\definecolor{test}{RGB}{100,100,100} % Couleurs pour les titres
\definecolor{test2}{RGB}{200,200,200} % Couleurs pour les titres



\usecolortheme[named=EricTemplate]{structure}
\setbeamercolor{section in toc}{fg=test} % TOC sections
\setbeamercolor{title}{fg=EricTemplate,bg=} 
\setbeamercolor{block title}{fg=white,bg=test}
\setbeamercolor{block body}{fg=black,bg=test2}


\setbeamertemplate{caption}[numbered]
\setbeamertemplate{blocks}[rounded][shadow=true]


\usepackage[utf8]{inputenc}
\usepackage[T1]{fontenc}    % use 8-bit T1 fonts
\usepackage{url}            % simple URL typesetting
\usepackage{booktabs}       % professional-quality tables
\usepackage{amsfonts}       % blackboard math symbols
\usepackage{nicefrac}       % compact symbols for 1/2, etc.
\usepackage{microtype}      % microtypography
\usepackage{mathtools}
\usepackage{multicol}
\usepackage{comment}
\usepackage{forest}
\usepackage[linesnumbered, ruled]{algorithm2e}
\usepackage{algorithmic}
\usepackage{xcolor}
\usepackage{colortbl}
\usepackage{graphicx}
\usepackage{diagbox}
\usepackage{tikz}
\usetikzlibrary{shadows}
\usepackage{subcaption}
% \captionsetup[subfigure]{labelsep=space}
\usepackage{amsmath}
\usepackage{dsfont}
\usepackage{stmaryrd}
\usepackage[style=authoryear-comp,backend=bibtex]{biblatex}
%\bibliographystyle{apalike}
\addbibresource{thesiscite.bib}

\renewcommand{\footnotesize}{\tiny}
\newcommand{\mathbi}[1]{\boldsymbol{#1}}
\newtheorem{proposition}{Proposition}

%------------------------------------------------------------
%This block of code defines the information to appear in the
%Title page
\title[New template]{}

\author[J. Jacques]{}


\institute[ERIC]{} % (optional)

\date[January 2024]{} % (optional)


%End of title page configuration block
%------------------------------------------------------------



%------------------------------------------------------------
%The next block of commands puts the table of contents at the 
%beginning of each section and highlights the current section:

\AtBeginSection[]
{
  \begin{frame}
    \frametitle{Table of Contents}
    \tableofcontents[currentsection]
  \end{frame}
}
%------------------------------------------------------------

\setbeamertemplate{title page}
{
  \vbox{}
  \vfill
  \begin{centering}
    \begin{beamercolorbox}[sep=8pt,centerw]{title}
      \usebeamerfont{title}\inserttitle
    \end{beamercolorbox}
    \setbeamercolor{title}{bg=white,fg=structure}
    \begin{beamercolorbox}[sep=8pt,center]{title}
      {\usebeamercolor[fg]{titlegraphic}\inserttitlegraphic\par}
      \ifx\insertsubtitle\@empty%
      \else%
        \vskip0.25em%
        {\usebeamerfont{subtitle}\usebeamercolor[fg]{subtitle}\insertsubtitle\par}%
      \fi%     
    \end{beamercolorbox}%
    \vskip1em\par
    \begin{beamercolorbox}[sep=8pt,center]{author}
      \usebeamerfont{author}\insertauthor
    \end{beamercolorbox}
    \vskip-1em\par % change here
    \begin{beamercolorbox}[sep=8pt,center]{institute}
      \usebeamerfont{institute}\insertinstitute
    \end{beamercolorbox}
    \begin{beamercolorbox}[sep=8pt,center]{date}
      \usebeamerfont{date}\insertdate
    \end{beamercolorbox}\vskip0.5em
  \end{centering}
  \vfill
}


\begin{document}

% The next statement creates the title page.
\setbeamertemplate{background canvas}{
	\includegraphics[width=\paperwidth,height=\paperheight]{images/bg_slide_1.png}
    
}
\begin{frame}[plain]
\begin{tikzpicture}[remember picture,overlay]
\node
[   draw,
    drop shadow,
    above=0.5cm,
    align=center,
    fill=orange!100,
    rounded corners,
    inner xsep=15pt,
    inner ysep=10pt, 
    minimum width=0.65\textwidth,
    text width=0.6\textwidth
] (title) at (current page.center)
{
    \LARGE My Title  \\[5pt]
    \small My subtitle
};
\node
[   draw,
    drop shadow,
    below=0.5cm,
    align=center,
    fill=gray!20,
    rounded corners,
    inner xsep=15pt,
    inner ysep=10pt, 
    minimum width=0.4\textwidth,
    text width=0.35\textwidth
] (title) at (current page.center)
{
    \small My name  \\[3pt]
    \small My advisor\\
    \small Institute
};
\end{tikzpicture}


\end{frame}


% Set background for all pages
\setbeamertemplate{background canvas}{
	\includegraphics[width=\paperwidth,height=\paperheight]{images/bg_all.jpg}
}

\begin{frame}
\vspace{1cm}
Une phrase pour l'exemple
\end{frame}

%---------------------------------------------------------
%This block of code is for the table of contents after
%the title page
\begin{frame}
\frametitle{Table of Contents}
\tableofcontents
\end{frame}

%---------------------------------------------------------
\section{Introduction}
%---------------------------------------------------------

\begin{frame}
\frametitle{Motivation}

\end{frame}

\begin{frame}
\frametitle{Goal \& Scope}
\textbf{Goal:} Overall Goal
\begin{enumerate}
    \item Goal 1
    \item Goal 2
    \item Goal 3
\end{enumerate}
\vspace{0.5cm}
\pause
\textbf{Scope:}
\begin{itemize}
    \item Scope A
    \item Scope B
    \item Scope C
\end{itemize}
\end{frame}

% %---------------------------------------------------------
\section{Preliminaries and Problem Statement}
% %---------------------------------------------------------
\begin{frame}
\frametitle{Preliminaries}
\begin{block}{LE bloc} % Pour changer les couleurs, aller sur setbeamercolor{block ...}
    Un bloc
\end{block}
\end{frame}

\begin{frame}
\frametitle{Problem Statement}

\end{frame}

%---------------------------------------------------------
\section{Related Works}
%---------------------------------------------------------
\begin{frame}
\frametitle{Work 1}

\end{frame}

\begin{frame}
\frametitle{Work 2}

\end{frame}

\begin{frame}
\frametitle{Work 3}

\end{frame}

%---------------------------------------------------------
\section{Proposed Method}
%---------------------------------------------------------

\begin{frame}
\frametitle{Overview}

\end{frame}


\begin{frame}
\frametitle{Component 1}

\end{frame}

\begin{frame}
\frametitle{Component 2}

\end{frame}

\begin{frame}
\frametitle{Component 3}

\end{frame}

%---------------------------------------------------------
\section{Experiments and Results}
%---------------------------------------------------------
\begin{frame}
\frametitle{Datasets}

\end{frame}

\begin{frame}
\frametitle{Baselines}


\end{frame}

\begin{frame}
\frametitle{Metrics}

\end{frame}

\begin{frame}
\frametitle{First Result}

\end{frame}

\begin{frame}
\frametitle{Second Result}

\end{frame}

%---------------------------------------------------------
\section{Conclusion}
%---------------------------------------------------------
% Take away
% Contact
% Thank
\begin{frame}{Summary}

\end{frame}


\begin{frame}{Future Developments}

\end{frame}


\setbeamertemplate{footline}{}
\begin{frame}
    \centering
    \vspace{1cm}
    \textbf{\Huge{- THE END - }}\\
    \vspace{0.2cm}
    \textit{\Large{Thank you for your attention}}\\
    \vspace{15pt}
    Acknowledgements:\\
    I want to...\\
    \vspace{20pt}
    Contact:\\
    youremail@hcmut.edu.vn
\end{frame}

% \begin{frame}[allowframebreaks]
%         \frametitle{References}
%         \bibliographystyle{ieeetr}
%         \bibliography{reference}
% \end{frame}

\end{document}